\exercisesection{DIMENSION OF ALGEBRAS}

\begin{enumerate}
\def\labelenumi{\arabic{enumi})}
\item
  Let V be a discrete valuation ring, with fraction field K, and let n be an integer. Let A be the subring of the polynomial ring \(K[X_{1},...,X_{n}]\) formed by the polynomials whose constant term lies in V. The canonical homomorphism from V to A satisfies property (PM) and the canonical homomorphism \(V/m_{V} \rightarrow A/m_{V}A\) is an isomorphism. Moreover, one has \(\dim(A) \geqslant n + 1\) and \(\dim(V) = 1\) . As soon as \(n \geqslant 1\) , one therefore has \(\dim(A) > \dim(V) + \dim(A/m_{V}A)\) . Show that A is not Noetherian. Using formal series, construct an analogous example where A is local and not Noetherian.
\item
  Let V be a discrete valuation ring, with fraction field K and residue field k. Denote by A the ring \(K \times k\) and \(\rho : V \to A\) the homomorphism deduced from the canonical homomorphisms from V to K and k. The homomorphism \(\rho\) is injective, the map \(^{a}\rho : \text{Spec}(A) \to \text{Spec}(V)\) is surjective, A is a finitely generated V-algebra, and we have \(\dim(A) < \dim(V)\) .
\item
  Let A be an integral domain of dimension 1, and let \(\mathfrak{p}_0 \subset \mathfrak{p}_1 \subset \mathfrak{p}_2\) be three distinct nonzero prime ideals of A{[}X{]}. Show that we have \(\mathfrak{p}_0 \cap \mathbf{A} = \{0\}\) , \(\mathfrak{p}_1 = (\mathfrak{p}_1 \cap \mathbf{A})\cdot\mathrm{A}{[}X{]}\), and that \(\mathfrak{p}_2\) is a maximal ideal of A{[}X{]}.
\item
  Let A be a local integrally closed ring, \(m_A\) its maximal ideal and K its field of fractions. For every element x of K, we denote by A{[}x{]} the sub-A-algebra of K generated by x. Let x be an element of K* such that \(x^{-1}\) does not belong to A. Let A{[}X{]} be the ring of polynomials in one indeterminate X, with coefficients in A, and \(\varphi\) the evaluation map P \(\mapsto\) P(x) of A{[}X{]} on A{[}x{]}.
\end{enumerate}

\begin{enumerate}
\def\labelenumi{\alph{enumi})}
\item
  The kernel of \(\varphi\) is the ideal Q generated by the polynomials of the form \(aX + b\) , where \(a, b \in A\) and \(ax + b = 0\) . (If \(\mathrm{P}(X) = \sum_{i=0}^{n} a_i X^i\) is such that \(\mathrm{P}(x) = 0\) , note that then \(-b_0 = a_0 x\) is integral over \(A\) therefore belongs to \(A\) .)
\item
  The ideal \(m_A A[x]\) of \(A[x]\) is prime, and it is maximal if and only if \(x\) belongs to \(A\) .
\end{enumerate}

\begin{enumerate}
\def\labelenumi{\arabic{enumi})}
\setcounter{enumi}{4}
\item
  Let A be an integrally closed ring with fraction field K. Given a prime ideal p of A, and an element x of K such that \(x \notin A_{p}\) , and \(x^{-1} \notin A_{p}\) the ideal pA{[}x{]} of A{[}x{]} is prime; we have: pA{[}x{]} ∩ A = p and the canonical homomorphism (A/p){[}X{]} → A{[}x{]}/pA{[}x{]} is an isomorphism.
\item
  \begin{enumerate}
  \def\labelenumii{\alph{enumii})}
  \tightlist
  \item
    Using the previous exercises, show that if A is an integral domain of dimension 1, then dim(A{[}X{]}) = 2 if and only if all the localizations at a prime ideal of the integral closure of A are valuation rings.
  \end{enumerate}
\end{enumerate}

\begin{enumerate}
\def\labelenumi{\alph{enumi})}
\setcounter{enumi}{1}
\tightlist
\item
  Show that if one has \(\dim(\mathbf{A}) = n\) and \(\dim(\mathbf{A}[\mathbf{X}]) > n + 1\), there exists a prime ideal \(\mathfrak{p}\) of A such that \(\operatorname{ht}(\mathfrak{p}) = 1\) and either \(\dim(\mathbf{A}_{\mathfrak{p}}[\mathbf{X}]) > 2\) or \(\dim((\mathbf{A}/\mathfrak{p})[\mathbf{X}]) > \dim(\mathbf{A}/\mathfrak{p}) + 1\) .
\item
  Deduce that if A is Prüferian (VII, § 2, exerc. 12), then we have \(\dim(\mathbf{A}[\mathbf{X}_1, ..., \mathbf{X}_n]) = \dim(\mathbf{A}) + n\) , for every integer \(n \geqslant 0\) .
\end{enumerate}

\begin{enumerate}
\def\labelenumi{\arabic{enumi})}
\setcounter{enumi}{6}
\item
  Let B be an integrally closed ring, with field of fractions L. We set \(d = \dim(B)\) and \(t = \dim(B[X])\) . Let L' be a non-trivial extension of L in which L is algebraically closed, and let V be a discrete valuation ring (which is not a field) having L' as its residue field. We denote by K the field of fractions of V and by A the set of elements of V whose image in L' under the canonical homomorphism V → L' belongs to B. The ring A is integrally closed with field of fractions K. Let p denote the kernel of the surjective homomorphism A → B. Show that every non-zero prime ideal of A contains p, whence dim(A) = dim(B) + 1. Considering an element x of V whose image in L' does not belong to L, show that the ring of fractions A \(_{p}\) is not a valuation ring, and that pA \(_{p}\) {[}x{]} is not minimal among the non-zero prime ideals of A \(_{p}\) {[}x{]}. Using exercise 6, deduce that one has dim(A{[}x{]}) ≥ t + 2. Conclude that one has dim(A{[}x{]}) = t + 2 by a direct argument. Finally show that for every pair (d, t) of integers with d + 1 ≤ t ≤ 2d + 1, there exists an integral domain A of dimension d such that dim(A{[}X{]}) = t.
\item
  Show that the polynomial ring A{[}X{]} is a faithfully flat A-module, that for every maximal ideal m of A one has dim(A{[}X{]}/mA{[}X{]}) = 1, and deduce from the preceding exercise the existence of ring homomorphisms \(\rho: A \to B\) satisfying condition (PM) and such that \(\dim(B) > \dim(A) + \inf_{\mathfrak{m}\in S}\dim(B/\mathfrak{m}B)\) (where S is the set of maximal ideals of A).
\item
  Let A be an integral domain and B an integral domain containing A and integral over A. Let p be a prime ideal of A such that the integral closure of the ring of fractions \(A_{p}\) is a local ring. For every prime ideal q of B lying over p, one has ht(q) = ht(p) (use V, § 2, n° 1, prop. 2).
\end{enumerate}

¶ 10) a) Let A be an integral domain, K its field of fractions, and m a positive integer. Let \(t_1, \ldots, t_m\) be elements of K and \(\varphi: \mathrm{A}[\mathrm{X}_1, \ldots, \mathrm{X}_m] \to \mathrm{A}[t_1, \ldots, t_m]\) the unique homomorphism of A-algebras such that \(\varphi(\mathrm{X}_i) = t_i\) . Show that the height of the kernel ideal of \(\varphi\) is equal to m. b) Deduce that for every \(m\) -tuple ( \(t_1, \ldots, t_m\) ) of elements of K, one has \(\dim(\mathrm{A}[t_1, \ldots, t_m]) \leqslant \dim(\mathrm{A}[\mathrm{X}_1, \ldots, \mathrm{X}_m]) - m\) .

¶ 11) Let A be an integral domain and K its field of fractions. Let \((0) \subset \mathfrak{p}_1 \subset \ldots \subset \mathfrak{p}_h\) be a chain of prime ideals of A. Show that there exists a valuation ring V of K and a chain of prime ideals \((0) \subset \mathfrak{q}_1 \subset \ldots \subset \mathfrak{q}_h\) of V such that \(\mathfrak{q}_i \cap A = \mathfrak{p}_i\) . (Proceed by induction on \(h\) , reducing to the case where A is a local ring and using, for the case \(h = 1\) , exercise 7 of VI, § 1.)

\begin{enumerate}
\def\labelenumi{\arabic{enumi})}
\setcounter{enumi}{11}
\tightlist
\item
  \begin{enumerate}
  \def\labelenumii{\alph{enumii})}
  \tightlist
  \item
    Let A and K be as above, let B be an integral domain containing A and integral over A, and let L be the fraction field of B. Let \(k \geqslant 0\) be an integer. If, for every \(m\)-tuple \((t_1, \ldots, t_m)\) of elements of L, we have \(\dim(A[t_1, \ldots, t_m]) \leqslant k\), then, for every \(m\)-tuple \((u_1, \ldots, u_m)\) of elements of L, we have \(\dim(B[u_1, \ldots, u_m]) \geqslant k\). b) Let p be a prime ideal of A. Show that if, for all \(m\)-tuples \((t_1, \ldots, t_m)\) of elements of K, one has \(\dim(A[t_1, \ldots, t_m]) \leqslant k\), then we have for all \(m\)-tuples \((s_1, \ldots, s_m)\) of elements of the field of fractions of A/p, the inequality \(\dim((A/p)[s_1, \ldots, s_m]) \leqslant k - \text{ht}(p)\) .
  \end{enumerate}
\end{enumerate}

¶ 13) Let A be an integral domain and K its field of fractions. For every integer \(k \geqslant 0\) , the following properties are equivalent:

\begin{enumerate}
\def\labelenumi{(\roman{enumi})}
\item
  Every subring B of K containing A has dimension \(\leqslant k\) .
\item
  Every valuation ring V of K containing A has dimension \(\leqslant k\) .
\item
  For every choice of \(k\) elements \(t_1, \ldots, t_k\) of \(\mathbf{K}\) , one has \(\dim(A[t_1, \ldots, t_k]) \leqslant k\) .
\item
  We have \(\dim(A[X_1, ..., X_m]) \leqslant k + m\) for every integer \(m \geqslant 0\) .
\end{enumerate}

\begin{enumerate}
\def\labelenumi{\arabic{enumi})}
\setcounter{enumi}{13}
\tightlist
\item
  Let A be an integral domain, with field of fractions K, and \(m \in N\). Show that we have
\end{enumerate}

\[
\dim (\mathrm{A} [ \mathrm{X} _ {1},..., \mathrm{X} _ {m} ]) = m + \sup \left\{\dim (\mathrm{A} [ t _ {1},..., t _ {m} ]) | t _ {i} \in \mathrm{K}, 1 \leqslant i \leqslant m \right\}.
\]

(We will use Exercise 10). More generally, if F is a field, an extension of K, show that we have

\[
\dim \left(\mathrm{A} \left[ \mathrm{X} _ {1}, \dots , \mathrm{X} _ {m} \right]\right) = \sup \left\{\dim (\mathrm{A} \left[ t _ {1}, \dots , t _ {m} \right]) \mid t _ {i} \in \mathrm{F}, 1 \leqslant i \leqslant m \right\} + \sup (0, m - \deg . \operatorname{tr} _ {\mathbf {K}} (\mathrm{F})).
\]

\begin{enumerate}
\def\labelenumi{\arabic{enumi})}
\setcounter{enumi}{14}
\tightlist
\item
  With the notations of the previous exercises, show that if for some integer \(k \geqslant 0\) , one has \(\dim(A[X_1, ..., X_m]) \geqslant k + m + 1\) for a suitable integer \(m\), one has the inequality \(\dim(A[X_1, ..., X_k]) \geqslant 2k + 1\) .
\end{enumerate}

  ¶ 16) Let A be an integral domain and K its field of fractions. One calls the valuative dimension of A and denotes it by \(\dim_{\mathrm v}(\mathrm A)\) the real number \(\sup_{A \subset V \subset K} \dim(V)\) , where V ranges over the set of valuation rings of K containing A

\begin{enumerate}
\def\labelenumi{\alph{enumi})}
\item
  Show that one has \(\dim(\mathrm A) \leqslant \dim_{\mathrm v}(\mathrm A)\) (exercise 11).
\item
  If A is Prüferian (VII, § 2, exerc. 12), show that one has \(\dim(\mathrm A) = \dim_{\mathrm v}(\mathrm A)\).
\item
  Show that, for all \(m \in \mathbf{N}\) , one has
\end{enumerate}

\[
\dim_ {\mathbf {v}} (\mathrm{A} [ \mathrm{X} _ {1},..., \mathrm{X} _ {m} ]) = \dim_ {\mathbf {v}} (\mathrm{A}) + m.
\]

(Note that a valuation ring is Prüferian, and one will use b) and exercise 6, p.~84.)

\begin{enumerate}
\def\labelenumi{\alph{enumi})}
\setcounter{enumi}{3}
\tightlist
\item
  Show that if one has \(\dim_{\mathrm v}(\mathrm A) = s < \infty\), one has, for every integer \(m \geqslant s - 1\) the equality
\end{enumerate}

\[
\dim (\mathrm{A} [ \mathrm{X} _ {1},..., \mathrm{X} _ {m} ]) = m + s.
\]

\begin{enumerate}
\def\labelenumi{\alph{enumi})}
\setcounter{enumi}{4}
\tightlist
\item
  Let \(s \in \mathbf{N}\) . Show that we have \(\dim_{\mathrm{v}}(\mathbf{A}) = s\) if and only if \(\dim(\mathbf{A}[\mathbf{X}_1, ..., \mathbf{X}_s]) = 2s\) (using \(d\) ) and Exercises 13 and 14).
\end{enumerate}

Deduce that, for every integral Noetherian ring A, one has \(\dim(\mathrm{A}) = \dim_{\mathrm{v}}(\mathrm{A})\).

\begin{enumerate}
\def\labelenumi{\arabic{enumi})}
\setcounter{enumi}{16}
\tightlist
\item
  Let A be an integral domain. Set \(n_k(A) = \dim(A[X_1, ..., X_k])\) for every integer \(k \geqslant 0\). Let D be the set of sequences of integers \((n_i)_{i \geqslant 0}\) such that there exists an integral domain A with \(n_i = n_i(A)\) for every integer \(i \geqslant 0\) .
\end{enumerate}

\begin{enumerate}
\def\labelenumi{\alph{enumi})}
\tightlist
\item
  Prove the inequalities
\end{enumerate}

\[
\sup \left(n _ {0} (\mathrm{A}) + i + 1, n _ {i} (\mathrm{A}) + 1\right) \leqslant n _ {i + 1} (\mathrm{A}) \leqslant \inf \left(2 n _ {i} (\mathrm{A}) + 1, 2 ^ {i + 1} \left(n _ {0} (\mathrm{A}) + 1\right) - 1\right).
\]

\begin{enumerate}
\def\labelenumi{\alph{enumi})}
\setcounter{enumi}{1}
\tightlist
\item
  Show that the only sequence of integers beginning with 1, 2 and belonging to D is the sequence \(n_{i}=i+1\) . (Let A be such that \(n_{0}(A)=1\) , \(n_{1}(A)=2\) . By exerc. 6, p.~84, the integral closure of A is Prüferian. Use exerc. 6, c), p.~84 to conclude.)\\
\item
  Prove the inequalities
\end{enumerate}

\[
n _ {0} (\mathrm{A}) + i \leqslant n _ {i} (\mathrm{A}) \leqslant (n _ {0} (\mathrm{A}) + 1) (i + 1) - 1.
\]

Deduce that if one has \(n_0(\mathbf{A}) = 0\) , then we have \(n_i(\mathbf{A}) = i\) for every \(i \in \mathbf{N}\) .

\begin{enumerate}
\def\labelenumi{\alph{enumi})}
\setcounter{enumi}{3}
\item
  Show that the sequence of differences \(n_{i+1}(\mathbf{A}) - n_i(\mathbf{A})\) is stationary.
\item
  We use the notations of exerc. 7, p.~84. Let δ be the transcendence degree of L' over L. Then:
\end{enumerate}

\begin{enumerate}
\def\labelenumi{(\roman{enumi})}
\item
  If \(\delta = 0\) , one has \(n_i(A) = n_i(B) + 1\) for every \(i\) .
\item
  If \(1 \leqslant \delta < \infty\) , one has \(n_i(A) = n_i(B) + i + 1\) for \(0 \leqslant i \leqslant \delta\) and \(n_i(A) = n_i(B) + \delta + 1\) for \(i \geqslant \delta\) .
\item
  If \(\delta = \infty\) , one has \(n_i(A) = n_i(B) + i + 1\) for every \(i\) .
\end{enumerate}

(For a characterization of the elements of D, one may consult T. Parker, Amer. J. Math., 97 (1975), pp. 308–311.)

* 18) Let A (resp. B) denote the algebra C \(\{z_{1}, z_{2}, z_{3}\}\) (resp. C \(\{u, v\}\) ) of convergent series in the variables \(z_{1}, z_{2}, z_{3}\) (resp. \(u, v\) ) and \(\varphi: A \to B\) the unique algebra homomorphism such that \(\varphi(z_{1}) = uv\) , \(\varphi(z_{2}) = v(e^{u} - 1)\) , \(\varphi(z_{3}) = v\) . Show that \(\varphi\) is injective, and that we have dim(A) = 3 and dim(B) = 2. *

